\documentclass[10pt,a4paper]{amsart}
\usepackage[margin=19mm]{geometry}
\usepackage{amsmath,amssymb,mathtools}
\usepackage[scr=rsfs,cal=euler]{mathalfa}
\usepackage{xfrac}
\usepackage[unicode,hidelinks,bookmarks=false]{hyperref}
\newcommand{\ie}{i.e.\ }
\newcommand{\cf}{cf.\ }
\newcommand{\resp}{respectively\ }
\newcommand{\Subsection}{\subsection}
\providecommand{\nobreakdash}{\nobreak\hbox{-}}
\setlength{\emergencystretch}{2em}
\multlinegap0pt
\usepackage{amssymb}
\usepackage{tikz}
\usepackage{float}
\usepackage[shortlabels]{enumitem}
\newtheorem{lemma}{Lemma}
\newtheorem{theorem}{Theorem}
\title{A Boolean polynomial operator for the Collatz $3n+1$ problem}
\author{Mario DeFranco}
\date{Source: arXiv:2608.25791v1 (CC BY 4.0)}
\begin{document}
\maketitle

\noindent\emph{Source and licence.} A Boolean polynomial operator for the Collatz $3n+1$ problem. Version arXiv:2608.25791v1 (CC BY 4.0), distributed by arXiv under the Creative Commons Attribution 4.0 International licence, http://creativecommons.org/licenses/by/4.0/, as recorded in the arXiv licence metadata for this exact version and shown on the abstract page https://arxiv.org/abs/2608.25791v1. Full licence notice: \texttt{licenses/arxiv-2608.25791-CC-BY-4.0.txt}.\par\smallskip

\noindent\textit{Editorial scope.} Counted unit: the article's theorem sum (source0.tex lines 602--608). The article's complete original proof of it is delivered with it (source0.tex lines 609--684, one \texttt{proof} environment of 76 lines). No mathematics is written, repaired or re-derived: the statement and the proof are the article's own lines, in the article's own notation, and no retained line is substituted or re-flowed.  Retained uncounted as the source-local context the proof needs: lines 572--577.  Nothing was omitted from any retained span, so the package carries no omission record.  No retained line contains a citation, so the wrapper carries no bibliography and no bibliography entry.  No retained line contains a figure environment, an image-inclusion command or drawing code, so no image is delivered and the package declares no asset dependency.  Editorial formatting only, in addition: an AMS \texttt{amsart} preamble that restates the article's own declarations and macros used by the retained lines, namely the environments \texttt{lemma}, \texttt{theorem} on separate counters, together with the standard packages those lines need.  The retained environments print in this document as Lemma 1, Theorem 1 in order of appearance, whereas the article prints the same items with the numbering of its own full document.  The complete native source of the article as distributed in the arXiv e-print for this exact version is bundled unchanged as \texttt{sources/208598/source0.tex}.
\begin{lemma} \label{l V}
Recall that $T_1(a,b)$ denotes the set of skip-1 sets $S$ in $[a,b]$ such that every block of $S$ has odd length.   Then 
\[
\sum_{S \in T_1(a,b)} y(S) =  V_1(y,a,b). 
\]
\end{lemma}

\begin{theorem} \label{t 1 2 sum}
\begin{align}
\sum_{S \in \Lambda(t)} y(S)=& \sum_{j=0}^{\lfloor \frac{t-3}{4}\rfloor}\sum_{i =-1}^{t-4j-4}P(y,1,i) V_1(y,i+2,i+3+4j) U(y,i+5+4j,t) \label{l 2 sum}\\
& + \sum_{i=0}^{\lfloor \frac{t-2}{4}\rfloor} V_1(y,1,4i+1) U(4i+3, t) \label{l 1 sum}
\end{align}
where we denote $P(1,-1) = P(1,0) = U(t+1,t) = 1$. 
\end{theorem}

\begin{proof} 


Suppose $1 \leq a < b <t$. 

%let $H(t,r,s)$ denote the set of subsets
%\[
%H(t,r,s) = \{ [s,t] \cup M \colon  M \subseteq [1,r]\}.
%\]
%denoting $[r,t]  = \emptyset$ for $r > t$. For $Z \in H(t,a,b)$,  let $\Lambda_i(t,a,b, Z) \subseteq \Lambda_i(t,a,b)$ denote the set of those $S$ such %that 
%\[
 %\Lambda_i(t,a,b, Z)  = \{S \in \Lambda_i(t,a,b) \colon  S \backslash [a,b] = Z \}.
%\]
%Let $\bar{\Lambda}_i(t,a,b)$ denote the set
%\[
%\bar{\Lambda}_i(t,a,b)=\{\mathrm{lead}(S,t) \colon S \in \Lambda_i(t,a,b) \}.
%\]
By construction we have
\begin{align*}
\Lambda_0(t,a,b) &= \emptyset \text{ if } b-a+1 \not\equiv 2 \pmod 4 \\ 
 \Lambda_1(t,a,b) &= \emptyset \text{ if } b-a+1 \not\equiv 1 \pmod 4. 
\end{align*}
Thus 
\begin{align}
\sum_{S \in \Lambda(t)} y(S) =& \sum_{a=2}^{t-1} \sum_{i=0}^{\lfloor \frac{t-1-a}{4} \rfloor}((\sum_{S \in \Lambda_1(t,a,a+4i)} y(S))+ (\sum_{S \in \Lambda_2(t,a-1,a+4i)} y(S)))\label{2 sum}\\ 
&+ \sum_{i=0}^{\lfloor \frac{t-2}{4} \rfloor}\sum_{S \in \Lambda_1(t,1,1+4i)} y(S). \label{1 sum}
\end{align}

Assume $b-a \equiv 0 \pmod 4$. Suppose $S \in \Lambda_1(t,a,b)$. Then $S$ is of the form 
\[
S = Z \cup \mathrm{lead}(S,t) \cup [b+2,t]
\]
where $Z$ may be any set 
\[
Z \subseteq [1,a-3].
\]
This follows from the construction of lead sequences: if $W \in T_1(a,b)$ and \[W\cup Z \cup [b+2,t] \in \Lambda_1(t,a,b),\] then we must have $a-1 \notin Z$ and $a-2 \notin Z$. 
Thus
\begin{equation} \label{SZ1}
\sum_{S \in \Lambda_1(t,a,b)} y(S) = \left(\sum_{Z \subseteq [1,a-3]} y(Z) \right )\left (\sum_{W \in T_1(a,b)} y(W) \right ) U(b+2, t) .
%&= \sum_{Z \in H(t,a,b)} P(y; 1,a-2) U(y; b+2,t)\sum_{W \in \bar{\Lambda}_1(t,a,b)} y(W) \\ 
%&= 
\end{equation}

Now
\[
\sum_{W \in T_1(a,b)} y(W) = V_1(y; a,b))
\]
by Lemma \ref{l V}. 
Thus \eqref{SZ1} is equal to 
\begin{equation} \label{PVU}
P(y; 1,a-3)V_1(y; a,b)U(y; b+2,t).
\end{equation}
This shows line \eqref{1 sum} evaluates to line \eqref{l 1 sum}. 

Next, for $a \geq 2$, suppose $S \in \Lambda_0(t,a-1,b)$. Then $S$ is of the form 
\[
S =  Z \cup \mathrm{lead}(S,t) \cup [b+2,t]
\]
where $Z$ may be any set 
\[
Z \subseteq [1,a-3].
\]

\begin{align} 
\sum_{S \in \Lambda_0(t,a-1,b)} y(S) =& \left (\sum_{Z \subseteq [1,a-3]} y(Z) \right ) \left (\sum_{W \in T_0(a-1,b)} y(W) \right ) U(y; b+2, t) \nonumber  \\
=& P(y; 1,a-3)y(a-1)\left(\sum_{W \in T_1(a,b)} y(W) \right)U(y; b+2, t)  \nonumber \\ 
= & P(y; 1,a-3)y(a-1)V_1(y; a,b)U(y; b+2,t).\label{SZ0}
\end{align}

Adding results \eqref{PVU} and \eqref{SZ0} yields 
\[
P(y; 1,a-3)V_1(y; a-1,b)U(y; b+2,t).
\]
This shows that line \eqref{2 sum} evaluates to line \eqref{l 2 sum}. This completes the proof. 
\end{proof}

\end{document}
